\section{Introduction}
\label{sec:intro}

Object detection systems have achieved remarkable performance on standard benchmarks; however, their reliability degrades significantly in the presence of out-of-distribution (OoD) inputs. A critical manifestation of this issue is hallucinated detections, where models assign high confidence to non-existent or irrelevant objects. This poses serious concerns in real-world deployments such as autonomous driving, surveillance, and medical imaging.

Recent work, including \textit{Revisiting Out-of-Distribution Detection in Real-time Object Detection: From Benchmark Pitfalls to a New Mitigation Paradigm}, has highlighted limitations in existing OoD detection benchmarks and emphasized the need for robust mitigation strategies. Most prior approaches rely on confidence re-scoring, energy-based methods, or training with curated OoD datasets. While effective to some extent, these methods often ignore an important factor: the intrinsic signal properties of input regions that trigger false detections.

Deep neural networks are known to exhibit spectral bias, favoring certain frequency components in input data. In the context of object detection, this leads to an over-reliance on high-frequency textures and fine-grained patterns, which may not correspond to semantically meaningful objects. Furthermore, small bounding boxes often lack sufficient contextual information, making them more prone to misclassification and inflated objectness scores.

In this work, we explore the hypothesis that frequency characteristics and object scale jointly influence hallucination behavior in object detectors. Through empirical analysis, we observe that:

\begin{itemize}
    \item High-frequency regions (e.g., textured backgrounds) frequently trigger false positives
    \item Small object proposals exhibit unstable and overconfident objectness scores
    \item These effects are amplified in OoD settings
\end{itemize}

Motivated by these observations, we propose a frequency-aware objectness calibration mechanism that operates at the detection level. Our approach introduces:

\begin{itemize}
    \item Frequency-based penalties using Laplacian energy
    \item Objectness-aware suppression to control overconfident predictions
    \item Size-aware regularization for small bounding boxes
\end{itemize}

Unlike prior methods, our framework is:

\begin{itemize}
    \item Model-agnostic (works with YOLO variants)
    \item Training-free (post-hoc calibration)
    \item Low-compute (suitable for resource-constrained settings)
\end{itemize}

To validate our approach, we conduct extensive experiments on both in-distribution (Pascal VOC) and OoD datasets constructed from Caltech-256. We evaluate performance using:

\begin{itemize}
    \item Hallucination Rate (HR) and false positive reduction
    \item Confidence distribution shifts
    \item Explainability metrics (Occlusion Sensitivity, Pointing Game)
    \item Faithfulness metrics (Insertion/Deletion AUC)
    \item Localization consistency via Grad-CAM
\end{itemize}

Our results show that frequency-aware calibration significantly reduces hallucinations while improving interpretability, providing both practical improvements and diagnostic insights into OoD behavior.
